
\section{第一章 Go语言入门认知}
\subsection{1.1 语言定位与核心特性}
\begin{itemize}
\item Go语言诞生背景与发展迭代
\item 极简语法设计编程思想
\item 静态类型编译型语言属性
\item 原生轻量级协程并发优势
\item 高效编译执行运行性能
\item 跨平台多系统编译适配
\item 垃圾自动回收内存管理
\item 云计算容器主流应用场景
\end{itemize}

\subsection{1.2 开发环境搭建配置}
\begin{itemize}
\item Go SDK开发工具包作用组成
\item 操作系统对应版本安装流程
\item 系统环境变量配置含义
\item GOPATH工作目录规则规范
\item Go Modules模块化机制作用
\item 命令行基础编译指令使用
\item 命令行运行程序执行方式
\item 主流IDE开发工具功能使用
\item 项目整体目录层级结构划分
\end{itemize}

\subsection{1.3 基础语法书写规范}
\begin{itemize}
\item 标识符合法字符组成范围
\item 标识符大小写命名区分
\item 标识符驼峰命名行业规范
\item 系统关键字分类与作用
\item 保留字使用限制约束
\item 单行注释编写格式用途
\item 多行块注释编写格式用途
\item 程序语句结束分隔规则
\item 大括号代码块边界划分
\item 包声明语句编写规范
\item 外部包导入引用语法格式
\end{itemize}

\section{第二章 变量、常量与数据类型}
\subsection{2.1 变量完整定义与使用}
\begin{itemize}
\item 变量内存存储本质含义
\item var关键字标准声明格式
\item 短变量:=快速声明写法
\item 变量声明无赋值默认规则
\item 变量初始化赋值操作
\item 多变量批量同时定义
\item 变量局部作用域生效范围
\item 变量全局作用域生效范围
\item 变量类型自动推导机制
\item 变量二次赋值数值修改
\end{itemize}

\subsection{2.2 基础内置数据类型}
\begin{itemize}
\item 有符号整型取值范围与字节
\item 无符号整型存储特性
\item 固定位数整型精准数值存储
\item float32单精度浮点小数类型
\item float64双精度高精度浮点
\item bool布尔真假逻辑判定类型
\item byte字节底层字符存储
\item rune Unicode字符编码存储
\item string字符串不可变类型
\item complex复数数值特殊类型
\end{itemize}

\subsection{2.3 复合聚合数据类型}
\begin{itemize}
\item 数组固定长度连续存储特点
\item 切片动态可变长度引用结构
\item map键值映射无序集合类型
\item struct结构体自定义组合类型
\item 指针变量内存地址引用类型
\item 函数类型封装功能引用类型
\end{itemize}

\subsection{2.4 常量定义与枚举生成}
\begin{itemize}
\item const常量声明固定格式
\item 常量不可修改核心特性
\item 常量声明同步赋值规则
\item 多常量批量组合定义
\item iota常量自动枚举生成器
\item 枚举自增数值变化规则
\end{itemize}

\subsection{2.5 数据类型强制转换}
\begin{itemize}
\item 不同基础类型转换语法格式
\item 整型浮点互相转换数值变化
\item 高精度低精度转换数据取舍
\item 字符与数值类型互相转换
\item 转换异常数值边界判定
\end{itemize}

\section{第三章 运算符完整体系}
\subsection{3.1 算术运算符}
\begin{itemize}
\item 加法运算数值计算规则
\item 减法运算数值计算规则
\item 乘法运算数值计算规则
\item 除法整数小数运算区分
\item 取模余数运算限制条件
\item 前置自增运算执行顺序
\item 后置自增运算执行顺序
\item 前置自减运算执行顺序
\item 后置自减运算执行顺序
\end{itemize}

\subsection{3.2 赋值与复合赋值运算符}
\begin{itemize}
\item 基础赋值数据覆盖逻辑
\item 加法复合赋值简写运算
\item 减法复合赋值简写运算
\item 乘法复合赋值简写运算
\item 除法复合赋值简写运算
\item 取模复合赋值简写运算
\end{itemize}

\subsection{3.3 关系比较运算符}
\begin{itemize}
\item 等于对比判断运算
\item 不等于对比判断运算
\item 大于数值大小比较
\item 小于数值大小比较
\item 大于等于边界比较判定
\item 小于等于边界比较判定
\item 关系运算布尔结果返回规则
\end{itemize}

\subsection{3.4 逻辑运算符}
\begin{itemize}
\item 逻辑与全真条件判定规则
\item 逻辑或一真条件判定规则
\item 逻辑非布尔值取反运算
\item 逻辑短路与截断执行特性
\item 逻辑短路或截断执行特性
\item 多条件逻辑组合嵌套判断
\end{itemize}

\subsection{3.5 位运算与特殊运算符}
\begin{itemize}
\item 按位与二进制运算规则
\item 按位或二进制运算规则
\item 按位异或二进制运算规则
\item 按位取反二进制运算规则
\item 二进制左移位移运算逻辑
\item 二进制右移位移运算逻辑
\item 取地址&运算符作用用法
\item 解引用*取值运算符作用
\item 运算符整体优先级排序
\item 运算符左右结合执行顺序
\end{itemize}

\section{第四章 程序流程控制结构}
\subsection{4.1 顺序执行结构}
\begin{itemize}
\item 自上而下逐行代码执行逻辑
\item 顺序语句先后执行约束规则
\end{itemize}

\subsection{4.2 分支判断结构}
\begin{itemize}
\item if单分支条件判断语句
\item if-else双分支二选一逻辑
\item else if多分支区间匹配判断
\item switch多条件分支匹配语句
\item switch表达式取值约束范围
\item case分支匹配执行规则
\item Go语言无case穿透特性
\item fallthrough主动穿透用法
\item default默认兜底分支作用
\item 分支语句多层嵌套使用方式
\end{itemize}

\subsection{4.3 循环遍历结构}
\begin{itemize}
\item for基础循环完整三段格式
\item for条件判断简化循环写法
\item 无限循环构成与编写方式
\item for-range遍历数组切片容器
\item for-range遍历map键值集合
\item for-range遍历字符串字符
\end{itemize}

\subsection{4.4 循环跳转与标签控制}
\begin{itemize}
\item break终止跳出当前单层循环
\item continue跳过单次循环迭代
\item 自定义标签标记循环层级
\item 标签配合break跳出多层循环
\item 单层循环业务逻辑编写
\item 双层循环嵌套运行执行顺序
\end{itemize}

\section{第五章 聚合容器数据类型}
\subsection{5.1 数组数组容器}
\begin{itemize}
\item 数组固定长度定义语法
\item 数组静态初始化赋值
\item 数组动态初始化创建
\item 下标从零开始索引规则
\item 通过下标存取修改元素
\item 数组整体遍历读取数据
\item 数组长度获取内置方法
\end{itemize}

\subsection{5.2 切片Slice动态容器}
\begin{itemize}
\item 切片引用底层数组原理
\item 数组截取生成切片方式
\item 直接声明创建空白切片
\item len切片实际元素长度
\item cap切片底层容量大小
\item append追加元素扩容机制
\item copy切片数据拷贝操作
\item 切片截取子切片写法
\end{itemize}

\subsection{5.3 Map键值映射容器}
\begin{itemize}
\item map键值结构定义格式
\item map初始化内存开辟
\item 键唯一值重复存储规则
\item 新增键值对赋值操作
\item 根据键查询对应数值
\item 删除指定键值元素
\item 遍历map无序键值数据
\item 键存在性判断校验方式
\end{itemize}

\subsection{5.4 Struct结构体自定义类型}
\begin{itemize}
\item 结构体类型格式定义编写
\item 结构体内部成员变量组成
\item 结构体实例创建初始化
\item 字面量快速赋值创建实例
\item 点运算符访问结构体成员
\item 结构体成员赋值修改数据
\end{itemize}

\section{第六章 函数编程体系}
\subsection{6.1 函数定义声明与调用}
\begin{itemize}
\item 函数代码封装复用设计思想
\item func关键字函数定义格式
\item 函数名命名规范约束
\item 函数代码块功能逻辑编写
\item 函数调用触发执行流程
\item main主函数程序唯一入口
\item 自定义函数调用语法规则
\end{itemize}

\subsection{6.2 函数参数传递特性}
\begin{itemize}
\item 固定位置参数定义写法
\item 同类型参数简写声明方式
\item ...可变参数不定长度传参
\item 值传递数据拷贝底层原理
\item 值传递无法修改原始变量
\end{itemize}

\subsection{6.3 函数返回值处理}
\begin{itemize}
\item 单个返回值定义与返回规则
\item return语句返回数据操作
\item 多返回值同步返回写法
\item 匿名返回值简写定义格式
\item 返回值接收赋值存储方式
\item 无返回值空函数特性
\end{itemize}

\subsection{6.4 函数进阶高级用法}
\begin{itemize}
\item 匿名函数无函数名定义方式
\item 匿名函数即时调用执行
\item 函数作为参数回调传递
\item 函数赋值变量引用调用
\item 递归函数自身嵌套调用
\item 递归递进执行过程
\item 递归回溯返回流程
\item 递归终止边界条件设置
\end{itemize}

\section{第七章 指针与内存引用机制}
\subsection{7.1 指针基础核心概念}
\begin{itemize}
\item 程序内存地址编号含义
\item 指针变量存储地址数值
\item 取地址运算符获取变量地址
\item 解引用运算符读取地址数据
\item 通过指针修改原始变量值
\item nil空指针定义与判定
\item 指针变量占用内存空间大小
\end{itemize}

\subsection{7.2 指针实际场景应用}
\begin{itemize}
\item 指针配合函数修改实参数据
\item 结构体指针快速访问成员
\item 指针数组存储多组地址
\item 多级指针嵌套地址引用
\end{itemize}

\section{第八章 包、模块与项目依赖管理}
\subsection{8.1 Package包基础管理}
\begin{itemize}
\item 包划分代码模块作用
\item 源文件首行包声明规范
\item 同目录同包文件访问规则
\item 大写开头公有成员外部访问
\item 小写开头私有成员同包访问
\item 内部包引用调用方式
\end{itemize}

\subsection{8.2 Module模块化工程}
\begin{itemize}
\item go mod init模块初始化命令
\item go.mod配置文件内容含义
\item 本地模块层级引用调用
\item 远程第三方包导入写法
\item 依赖包自动下载拉取
\item 依赖版本锁定与更新管理
\end{itemize}

\section{第九章 错误与异常处理机制}
\subsection{9.1 错误基础接口概念}
\begin{itemize}
\item error系统错误接口定义
\item 错误信息字符串存储形式
\item 函数返回错误值判定逻辑
\end{itemize}

\subsection{9.2 常规错误捕获处理}
\begin{itemize}
\item errors创建基础错误对象
\item 自定义错误信息构造生成
\item 接收函数返回错误变量
\item 错误非空判断业务处理
\end{itemize}

\subsection{9.3 恐慌宕机与恢复机制}
\begin{itemize}
\item panic主动触发程序恐慌
\item 恐慌触发程序中断特性
\item defer延迟执行语句声明规则
\item defer栈后进先出执行顺序
\item recover捕获恐慌异常函数
\item 捕获异常恢复程序运行
\end{itemize}

\section{第十章 原生并发编程}
\subsection{10.1 Goroutine轻量级协程}
\begin{itemize}
\item 协程与线程进程区别关系
\item 协程极小内存占用优势
\item go关键字启动协程语法
\item 主协程与子协程执行逻辑
\item 多协程无序并发执行特点
\end{itemize}

\subsection{10.2 Channel协程通信通道}
\begin{itemize}
\item 通道实现协程数据通信原理
\item 无缓冲通道同步收发特性
\item 有缓冲通道异步收发原理
\item 通道声明创建初始化
\item 通道发送数据写入操作
\item 通道接收数据读取操作
\item 通道关闭标识与判断
\end{itemize}

\subsection{10.3 并发同步控制}
\begin{itemize}
\item 多协程共享数据竞争问题
\item sync.WaitGroup等待协程全部结束
\item 等待组计数器增减规则
\item sync.Mutex互斥锁锁定临界代码
\item 加锁解锁规避数据错乱
\end{itemize}

\section{第十一章 接口与反射机制}
\subsection{11.1 Interface接口抽象类型}
\begin{itemize}
\item 接口方法集合定义规范
\item 结构体隐式实现接口特性
\item 实现全部方法即匹配接口
\item 接口变量存储结构体实例
\item 接口多态动态调用方法
\item 空接口interface{}通用存储
\item 空接口接收任意数据类型
\end{itemize}

\subsection{11.2 运行时反射机制}
\begin{itemize}
\item reflect包反射核心功能
\item 反射获取变量基础类型
\item 反射获取变量具体数值
\item 反射修改变量内部数据
\item 反射解析结构体成员信息
\item 反射动态调用绑定方法
\end{itemize}

\section{第十二章 文件IO与目录管理}
\subsection{12.1 基础文件读写操作}
\begin{itemize}
\item os包文件打开创建方式
\item 文件打开权限模式设置
\item 文件关闭资源释放规范
\item 字节数组批量读取文件
\item 单行文本内容读取操作
\item 数据覆盖写入文件内容
\item 追加模式写入新增数据
\end{itemize}

\subsection{12.2 目录与文件属性管理}
\begin{itemize}
\item 多级目录创建指令操作
\item 空目录与文件删除操作
\item 获取文件大小修改时间属性
\item 遍历指定目录全部文件
\item 文件重命名移动路径操作
\end{itemize}

\section{第十三章 常用标准库使用}
\subsection{13.1 通用基础工具库}
\begin{itemize}
\item strings字符串截取替换分割
\item strconv字符串与数值互转
\item time时间获取格式化计算
\item math数学常用运算函数
\item rand随机数生成范围控制
\end{itemize}

\subsection{13.2 简易网络编程}
\begin{itemize}
\item net包基础网络连接原理
\item TCP服务端端口监听绑定
\item 客户端发起TCP连接请求
\item 两端网络数据收发通信
\end{itemize}
